\documentclass{article}

\usepackage{booktabs}
\usepackage{tabularx}

\title{Development Plan\\\progname}

\author{\authname}

\date{}

\input{../Comments.text}
\input{../Common.text}

\begin{document}

\maketitle

\begin{table}[hp]
\caption{Revision History} \label{TblRevisionHistory}
\begin{tabularx}{\textwidth}{llX}
\toprule
\textbf{Date} & \textbf{Developer(s)} & \textbf{Change}\\
\midrule
Sep 24th, 2026 & All & Added first draft version\\
Date2 & Name(s) & Description of changes\\
... & ... & ...\\
\bottomrule
\end{tabularx}
\end{table}

\newpage{}

\wss{Put your introductory blurb here.  Often the blurb is a brief roadmap of
what is contained in the report.}

\wss{Additional information on the development plan can be found in the
\href{https://gitlab.cas.mcmaster.ca/courses/capstone/-/blob/main/Lectures/L02b_POCAndDevPlan/POCAndDevPlan.pdf?ref_type=heads}
{lecture slides}.}

\section{Confidential Information?}

\wss{State whether your project has confidential information from industry, or
not.  If there is confidential information, point to the agreement you have in
place.}

\wss{For most teams this section will just state that there is no confidential
information to protect.}
\section{IP to Protect}

\wss{State whether there is IP to protect.  If there is, point to the agreement.
All students who are working on a project that requires an IP agreement are also
required to sign the ``Intellectual Property Guide Acknowledgement.''}

\section{Copyright License}

\wss{What copyright license is your team adopting.  Point to the license in your
repo.}

\section{Team Meeting Plan}
Weekly meetings will be held between all members on Wednesday at 6:30PM, ideally
in person; whether that's on campus or at a team member's home. This
schedule is flexible to accomodate online meetings, if necessary. Members are expected to have
read up on all relevant materials and completed any assigned action items from the previous meeting,
to ensure that meets are productive and efficient. All members are encouraged to chair meetings, and
meeting minutes will be taken by a designated team member, and shared with the team after each meeting.

Additionally, weekly meetings with the supervisor and other professors involved in the project will be
held on Fridays, in order to stay on track with requirements and receive
ongoing guidance throughout our deliverables. Each meeting will be recorded and meeting minutes will
be taken as well with the same template as general meetings.
\\
\\
\underline{Meeting Minutes/Agenda Template}

\begin{enumerate}
    \item Date, Time, Location
    \item Attendees
    \item Goal of Meeting
    \item Agenda items (Review of previous meeting minutes, discussion of any challenges the team has, questions/concerns, and required team decisions)
    \item Action Items (Assign responsibilities to each member, note down deadlines and potential blockers between items)
    \item \textbf{IF a supervisor meeting is upcoming}: note down all questions to be discussed with the supervisor, and any updates that we can provide them.
    \item \textbf{IF a new deliverable was dropped}: review requirements as a team, brainstorm ideas/thoughts, and assign responsibilities to each member.
    \item Notes (Summarize decisions, close meeting)
\end{enumerate}


\section{Team Communication Plan}

\wss{Issues on GitHub should be part of your communication plan.}

\section{Team Member Roles}

\wss{You should identify the types of roles you anticipate, like notetaker,
leader, meeting chair, reviewer.  Assigning specific people to those roles is
not necessary at this stage.  In a student team the role of the individuals will
likely change throughout the year.}

\section{Workflow Plan}

\subsection*{Git Usage}
The team will use Git and GitHub for version control and project management. After being
assigned individual tasks, team members are required to create an issue, and to create
and work on a new branch. On completion of their work, they will submit
a pull request with a description of the changes made, and all steps to test said changes
to add the reviewer. All members are encouraged to review any open pull request, but 
at least one team member must review and approve the changes before the user
can merge their branch onto main. 

\subsection*{Issue Management}
All progress on the team will be tracked via Github Issues and a Kanban board. 
This includes every lecture and different type of meeting attended, as well as all
development tasks and deliverables. Template issues have been created for
each type of task/issue, and can be easily updated via configurable files in the repository;
currently including sections for dates, descriptions, and relevant notes. Each issue also
has a specific label to categorize it (e.g ``Lecture'' for Lecture Issues, and ``Bug''
all Development Bug Issues).
Each member will be responsible for updating their own issues with relevant pull requests,
and status changes as the task progresses. The coordinator will also regularly 
review the project board to ensure that all issues are on track.
\\\\
\underline{Issue Templates}:
\begin{enumerate}
    \item Lecture Issue: To track any relevant notes detailed in lecture.
    \item Supervisor Meeting Issue: To track all pre-meeting agenda items, questions, and notes with the supervisor and external teams.
    \item Team Meeting Issue: To track all pre-meeting agenda items, questions, and notes for general meetings.
    \item TA Meeting Issue: To track all pre-meeting agenda items, questions, and notes with the TA.
    \item Deliverable Issue: To track requirements and progress on all deliverable items.
    \item Blank Issue: To track any other important items.
\end{enumerate}

\subsection*{CI/CD}
To prevent any failures or issues from occurring, we'll be using Github Actions to implement
our main CI/CD pipelines. These workflows will then be triggered on every pull request
that opens or updates, and again on every merge to main. Each workflow will run 
the project's integration, contract or unit tests, and check for any formatting or
linting errors before allowing the changes to be merged. If any of these builds or runs
fail, the team member who opened the pull request must go back and fix the issue before 
merging. Another workflow will also be added to compile our LaTeX documentation whenever a
\texttt{.tex} file is changed, ensuring that every document still generates a PDF without
any errors. The compiled PDFs will also be uploaded as artifacts so whoever is reviewing
the pull request can actually verify the final formatted output before approving.
Once every check has passed, a final team member must stamp their approval in order to
finally merge the changes into main.
\\\\
\underline{Workflow Summary}:
\begin{enumerate}
    \item Brainstorm and discuss all tasks as a team.
    \item Individual task is assigned to a team member, who creates an issue on the kanban board to track.
    \item Team member creates a new branch from latest main, to work on their task. They must ensure that they're frequently pruning their branch to the latest main.
    \item Any progress or significant changes should be updated on the kanban issue board.
    \item Team member regularly commits their changes, with descriptive/clear commit messages.
    \item Once the task is complete, the team member submits a pull request, and assigns a reviewer. The pull request should include a description of the changes made, and any relevant testing information.
    \item The reviewer checks the PR, following any testing or tophatting instructions provided, and leaves comments if it needs any changes. If not, they approve the PR with a smiley face.
    \item Once the PR is approved, the team member double checks their work and merges their branch into main.
    \item Team member closes the issue on the kanban board.
\end{enumerate}

\section{Project Decomposition and Scheduling}

\begin{itemize}
  \item How will you be using GitHub projects?
  \item Include a link to your GitHub project
\end{itemize}

\wss{How will the project be scheduled?  This is the big picture schedule, not
details. You will need to reproduce information that is in the course outline
for deadlines.}

\section{Proof of Concept Demonstration Plan}

What is the main risk, or risks, for the success of your project?  What will you
demonstrate during your proof of concept demonstration to convince yourself that
you will be able to overcome this risk?

\section{Expected Technology}

\wss{What programming language or languages do you expect to use?  What external
libraries?  What frameworks?  What technologies.  Are there major components of
the implementation that you expect you will implement, despite the existence of
libraries that provide the required functionality.  For projects with machine
learning, will you use pre-trained models, or be training your own model?  }

\wss{The implementation decisions can, and likely will, change over the course
of the project.  The initial documentation should be written in an abstract way;
it should be agnostic of the implementation choices, unless the implementation
choices are project constraints.  However, recording our initial thoughts on
implementation helps understand the challenge level and feasibility of a
project.  It may also help with early identification of areas where project
members will need to augment their training.}

Topics to discuss include the following:

\begin{itemize}
\item Specific programming language
\item Specific libraries
\item Pre-trained models
\item Specific linter tool (if appropriate)
\item Specific unit testing framework
\item Investigation of code coverage measuring tools
\item Specific plans for Continuous Integration (CI), or an explanation that CI
  is not being done
\item Specific performance measuring tools (like Valgrind), if
  appropriate
\item Tools you will likely be using?
\end{itemize}

\wss{git, GitHub and GitHub projects should be part of your technology.}

\section{Use of AI and LLMs}

\wss{This section is a continuation of the previous section.  The role of LLMs
in your project as important enough to warrant their own section.  
How will your team be using LLMs for code and documentation development?  
What tools will you be using?  How will you verify the output of your LLMs?}

\section{Coding Standard}

\wss{What coding standard will you adopt?}

\newpage{}

\section{Appendix --- Generative AI Use Disclosure}

\wss{ChatGPT (version 5) was used to brainstorm ideas for the template for this disclose section.}

\wss{This section is about the use of AI to write this document, not your
planned use of AI for your project.  You discuss that topic in one of the
sections in the body of this report.}

\subsection{AI Tools Used}

\wss{Give the name of the AI tool(s) used and the version.}

\subsection{How and Where Used}

\wss{Explain what the tool(s) were used for.  Examples include: brainstorming,
explaining a technical concept, reviewing the document, generating or modifying
text, generating or modifying diagrams, identifying errors or omissions or
inconsistencies.}

\wss{For each of the uses, identify which parts of the document were affected.}

\wss{You do not need to report every interaction, but you should disclose the 
uses that materially contributed to your work}

\subsection{Verification of AI-Generated Content}

\wss{\begin{itemize}
    \item Describe how the team checked AI-generated or AI-assisted material
    before including it in the document.
    \item In particular, verify factual claims, technical assertions,
    requirements, calculations, references, and other information for which an
    AI system may produce incorrect or fabricated results.
    \item \textbf{The team is responsible for the accuracy, completeness,
consistency, and appropriateness of the submitted document, regardless of
whether material was generated by a human or an AI system.}
\end{itemize}}

\newpage{}

\section*{Appendix --- Reflection}

\wss{Not required for CAS 741}

\input{../Reflection.text}

\begin{enumerate}
    \item Why is it important to create a development plan prior to starting the
    project?
    \item In your opinion, what are the advantages and disadvantages of using
    CI/CD?
    \item What disagreements did your group have in this deliverable, if any,
    and how did you resolve them?
\end{enumerate}

\subsection*{Catherine}

Without a development plan to layout the foundation for how our project
is structured and run, I think we'd immediately lose our heads when it comes
to things like formatting, timelines, and most \textit{importantly} motivation. 
When things get messy and convoluted, I personally dread tackling it and putting
off work as important as our capstone project is the last thing I want to do.
The plan is there to make things easier for ourselves as members, as it's 
just a structured guideline that we can continuously refer to throughout the 
project. We don't have to continuously reinvent new rules or reconvene to 
discuss things that we've already done and gone over. As the project progresses,
I'm 100\% sure things will only become more complex and messy, so having a plan in 
place from the start is even more important to ensure that we're moving as efficiently
and as organized as possible.

I believe that there's a lot of pros to having CI/CD integrated in your project.
If there's any test-breaking code, linting errors that create rendering
or deployment errors, or stale code from bad merges/previous commits; 
having a pipeline to catch everything \textit{before} they go into main
can ultimately save a lot of headache everybody involved. As now maintenance can happen
while the new changes are contained in a pull request. On the downside though, they can take a while
to setup. I expect that we'll have to set aside additional time for the 
testing process of any additional pipelines we want to create.

Finally, I don't believe we had any team disagreements for this deliverable. We had a 
pretty straightforward and productive meeting on our next steps, brainstorming
collectively for each section item, before we individually assigned tasks and
broke off to work on them separately. We agreed though that if we have any future
disagreements, the members involved would communicate respectfully and sort it out 
together before moving forward with any work. 


\newpage{}

\section*{Appendix --- Team Charter}

\wss{borrows from
\href{https://engineering.up.edu/industry_partnerships/files/team-charter.pdf}
{University of Portland Team Charter}}

\subsection*{External Goals}

For this capstone project, our team's external goals are to achieve an A+ first and foremost,
but to also have a finished product that actually demonstrates what we've learned about
healthcare in a speech-language pathology setting, and also how technology can actually 
reduce barriers that many in the field face, from both clinicians and patients. We hope
to have a product that future employers can utilize, and also showcase on our resumes
to show off our skillset as engineers.

\subsection*{Attendance}

\subsubsection*{Expectations}

\wss{What are your team's expectations regarding meeting attendance (being on
time, leaving early, missing meetings, etc.)?}

\subsubsection*{Acceptable Excuse}

\wss{What constitutes an acceptable excuse for missing a meeting or a deadline?
What types of excuses will not be considered acceptable?}

\subsubsection*{In Case of Emergency}

\wss{What process will team members follow if they have an emergency and cannot
attend a team meeting or complete their individual work promised for a team
deliverable?}

\subsection*{Accountability and Teamwork}

\subsubsection*{Quality} 

\wss{What are your team's expectations regarding the quality
of team members' preparation for team meetings and the quality of the
deliverables that members bring to the team?}

\subsubsection*{Attitude}

\wss{What are your team's expectations regarding team members' ideas,
interactions with the team, cooperation, attitudes, and anything else regarding
team member contributions?  Do you want to introduce a code of conduct?  Do you
want a conflict resolution plan?  Can adopt existing codes of conduct.}

\subsubsection*{Stay on Track}

\wss{What methods will be used to keep the team on track? How will your team
ensure that members contribute as expected to the team and that the team
performs as expected? How will your team reward members who do well and manage
members whose performance is below expectations?  What are the consequences for
someone not contributing their fair share?}

\wss{You may wish to use the project management metrics collected for the TA and
instructor for this.}

\wss{You can set target metrics for attendance, commits, etc.  What are the
consequences if someone doesn't hit their targets?  Do they need to bring the
coffee to the next team meeting?  Does the team need to make an appointment with
their TA, or the instructor?  Are there incentives for reaching targets early?}

\subsubsection*{Team Building}

\wss{How will you build team cohesion (fun time, group rituals, etc.)? }

\subsubsection*{Decision Making} 

\wss{How will you make decisions in your group? Consensus?  Vote? How will you
handle disagreements? }

\end{document}
